\documentclass[11pt]{amsart}

\usepackage[T1]{fontenc}
\usepackage{lmodern}
\usepackage{microtype}
\usepackage{mathtools,amssymb,amsthm}
\usepackage[hidelinks]{hyperref}

\hypersetup{
  pdftitle={An identifying code of density 59/156 in the 6-row square strip}
}

\newtheorem{theorem}{Theorem}
\newtheorem{lemma}[theorem]{Lemma}
\newtheorem{proposition}[theorem]{Proposition}
\theoremstyle{remark}
\newtheorem*{remark}{Remark}

\newcommand{\Ncl}{N}
\newcommand{\dstar}{d^{*}}

\title[An identifying code of density 59/156 in the 6-row square strip]
{An Identifying Code of Density $59/156$ in the 6-Row Square Strip}

\date{}

\subjclass[2020]{05C69, 05C63, 94B25}
\keywords{identifying code, minimum density, infinite square strip, periodic code,
explicit construction}

\begin{document}

\begin{abstract}
Let $S_k$ be the subgraph of the square lattice induced by
$\mathbb Z\times\{0,\dots,k-1\}$ and let $\dstar(S_k)$ be the minimum density of an
identifying code of $S_k$ for radius $1$. Theorem~30 of the Lobstein--Hudry--Charon
survey records $\dstar(S_k)$ for $k\le5$ and, for $k\ge6$, the bracket
$7/20+1/(20k)\le \dstar(S_k)\le 7/20+3/(10k)$, which at $k=6$ is
$[43/120,\,2/5]$ of width $1/24$. We exhibit an explicit period-$52$ identifying code
of $S_6$ of density $59/156$, printed here in full, giving
\[
  \dstar(S_6)\ \le\ \frac{59}{156}=0.3782051282\ldots\ <\ \frac25 ,
\]
so the recorded interval at $k=6$ narrows from above; nothing in the source is contradicted.
Verifying the code from the definition is a finite check of $312$ vertices and $1612$ vertex
pairs and needs no solver; the accompanying program carries it out. Only the upper bound is
established here; the exact value of $\dstar(S_6)$ is left open, and no claim is made that
$59/156$ is close to it.
\end{abstract}

\maketitle

\section{The source statement}

Let $S_k$ denote the subgraph of the square lattice induced by
$\mathbb Z\times\{0,\dots,k-1\}$, the \emph{$k$-row infinite square strip}. For a vertex
$v$ write $\Ncl[v]$ for its \emph{closed} neighbourhood in $S_k$. A set
$C\subseteq V(S_k)$ is an \emph{identifying code} (of radius $1$) if
\begin{equation}\label{eq:idc}
  \Ncl[u]\cap C\neq\emptyset \quad\text{for every }u,
  \qquad
  \Ncl[u]\cap C\neq \Ncl[v]\cap C \quad\text{for all } u\neq v .
\end{equation}
The set $\Ncl[u]\cap C$ is the \emph{identifier} of $u$. We take the \emph{density} of $C$ and
the invariant $\dstar$ from the source, so that every statement below is about exactly the
quantity Theorem~30 brackets: \cite[Section~6.4.1]{LHC} defines the density of a code $C$ in
$G_\pi^{[k]}$, for every $k\ge1$, as
$\limsup_{n\to\infty}|C\cap N_n[v]|/|N_n[v]|$, where $N_n[v]$ is the ball of radius $n$ about an
arbitrary vertex $v$, and $\dstar(S_k)=\partial_1^{Id}(G_S^{[k]})$ is the infimum of the densities
of the identifying codes of $S_k$.

The code exhibited in Section~2 is periodic with period $52$, so its density is a limit rather than
merely a $\limsup$, and it takes the same value if the balls $N_n[v]$ are replaced by rectangles
$[-n,n]\times\{0,\dots,k-1\}$: at $k=6$ the ball $N_n[v]$ and the rectangle
$[-n,n]\times\{0,\dots,5\}$ agree except on at most $6\cdot6=36$ cells near the two ends, out of
the $6(2n+1)$ cells of the rectangle, so the two window ratios differ by $O(1/n)$ and have the
same limit. Either window therefore yields the density computed in Section~\ref{sec:bound}, which
the program evaluates as the exact rational $118/312$ over one period.

The state of knowledge is recorded in A.~Lobstein, O.~Hudry and I.~Charon,
\emph{Locating-Domination and Identification}, Chapter~6 of \emph{Topics in Domination in
Graphs} \cite{LHC}, Section~6.4.1.1 ``The Square Grid'', where $\partial_1^{Id}$ is that
chapter's symbol for the minimum density we call $\dstar$ and the superscript $[k]$ denotes
the strip of height $k$. Its Theorem~30 reads, in part (with the source's stacked fractions
written inline and parentheses added in $1/(20k)$ and $3/(10k)$):

\begin{quote}
\textbf{Theorem 30}
(f) ([23]) For $k\ge6$, we have
\[
  7/20+1/(20k)\le \partial_1^{Id}(G_S^{[k]})\le 7/20+3/(10k).
\]
\end{quote}

Parts (a)--(e) record exact values for $k=1,\dots,5$; for $k\ge6$ the source records only the
bracket of part~(f). At $k=6$ that bracket is
\[
  \Bigl[\tfrac{7}{20}+\tfrac{1}{120},\ \tfrac{7}{20}+\tfrac{3}{60}\Bigr]
  = \Bigl[\tfrac{43}{120},\ \tfrac{2}{5}\Bigr],
  \qquad \tfrac{2}{5}-\tfrac{43}{120}=\tfrac{1}{24},
\]
the upper endpoint being exactly $2/5$. Part~(f) is a proved theorem carrying a gap at every
$k\ge6$; the source leaves the value at $k=6$ undetermined, recording only the interval
$[43/120,2/5]$. What the present note does is narrow that recorded interval from above.

\medskip
\noindent\textbf{Attribution.} The bracket endpoints of Theorem~30(f) are due to Bouznif,
Havet and Preissmann \cite{BHP}. Since the transfer-matrix and periodicity machinery of
\cite{Jiang}, \cite{BHP} and \cite{BDMP} could in principle evaluate $k=6$, we record what each
of the works we read reaches: \cite{DGM} gives $k=1,2$ and the first $k\ge6$ bracket; \cite{BHP}
proves the two bracket endpoints and, as an exact value, only $\dstar(S_3)=7/18$; \cite{Jiang}
supplies the periodicity theorem and the minimum-cycle-mean route and evaluates $k=4$ and $k=5$;
\cite{BDMP} resolves heights $1,2,3,4$. None of them states a value or a code at $k=6$, which is
why \cite{LHC} still records that cell as an interval in 2020. No priority claim is made beyond
that: no search was carried out past the works listed here, and $59/156$ or a better bound may
well be recorded elsewhere. What is offered is the code of Section~2 and the verification of
\eqref{eq:idc} for it.

\section{The object, and the upper bound}\label{sec:bound}

Let $L=52$, let the $52$ integers below be $c_0,\dots,c_{51}$, read left to right and top
to bottom, and put
\[
  C=\bigl\{(x,y)\in\mathbb Z\times\{0,\dots,5\}\ :\ \text{bit } y \text{ of } c_{x \bmod 52}
  \text{ equals } 1\bigr\}.
\]

\begin{center}
\footnotesize
\begin{tabular}{l}
\texttt{18 24 19 2 58 2 41 10 40 3 58 0 21 21 21 0 58 3}\\
\texttt{40 10 41 2 58 2 22 16 51 2 26 16 23 16 37 20 5 48}\\
\texttt{23 0 42 41 42 0 23 48 5 20 37 16 23 16 50 6}
\end{tabular}
\end{center}

\noindent Every $c_x$ satisfies $0\le c_x<64$. Equivalently, writing \texttt{1} for a cell
in $C$ and a dot for a cell outside it, one period of $C$ is the following $6\times52$
array, row $y$ being the row $y$ of the strip:

\begin{center}
\footnotesize
\begin{tabular}{ll}
\texttt{row 0} & \texttt{..1...1..1..111..1..1.....1...1.1.1.1..1..1.1.1.1...}\\
\texttt{row 1} & \texttt{1.1111.1.11.....11.1.1111.111.1.....1.1.1.1.....1.11}\\
\texttt{row 2} & \texttt{............111.........1.....1.111.1.....1.111.1..1}\\
\texttt{row 3} & \texttt{.1..1.111.1.....1.111.1.....1.........111...........}\\
\texttt{row 4} & \texttt{111.1.....1.111.1.....1.111.1111.1.11.....11.1.1111.}\\
\texttt{row 5} & \texttt{....1.1.1.1.....1.1.1.1...1.....1..1..111..1..1...1.}
\end{tabular}
\end{center}

The row weights are $18$, $26$, $15$, $15$, $27$, $17$; they sum to $118$ code cells out of
the $6\cdot52=312$ cells of the fundamental domain.

\begin{theorem}\label{thm:upper}
The set $C$ above is an identifying code of $S_6$, of density $118/312=59/156$, so
\[
  \dstar(S_6)\le\frac{59}{156}<\frac{2}{5},
  \qquad \frac{2}{5}-\frac{59}{156}=\frac{17}{780}.
\]
\end{theorem}

\begin{proof}
$C$ is invariant under $(x,y)\mapsto(x+52,y)$ and meets a fundamental domain in $118$ of
$312$ cells, so its density is $118/312=59/156$. It remains to verify \eqref{eq:idc}, and
the point is that this is a \emph{finite} check with an explicit bound.

If $\Ncl[u]\cap \Ncl[v]\neq\emptyset$ then some vertex is at distance at most $1$ from
both $u$ and $v$, so $d(u,v)\le2$. Contrapositively, if $d(u,v)\ge3$ the two closed
neighbourhoods are disjoint; if moreover both identifiers are nonempty --- which is the
domination clause --- then they are distinct sets, since a nonempty set cannot equal a set
disjoint from it. So the separation clause of \eqref{eq:idc} only has to be verified for
pairs at distance at most $2$, and by the $52$-periodicity of $C$ only for pairs whose
left-most vertex lies in $[0,51]\times\{0,\dots,5\}$.

That is $312$ domination tests and $1612$ pair tests, each one an inspection of at most
$10$ cells of the printed array; all of them pass. The accompanying program carries out this
verification and reports no failures.
\end{proof}

\section{Scope}

$59/156$ lies strictly inside $[43/120,2/5]$, with $2/5-59/156=17/780$ and
$59/156-43/120=31/1560$: Theorem~\ref{thm:upper} narrows from above the interval the source
records at $k=6$; nothing in Theorem~30(f) is contradicted. The value of $\dstar(S_6)$, and
the cases $k\ge7$, are left open here.

The tile of Section~2 was found by a randomised local search over periodic patterns of period
dividing $52$, and $52$ was simply the period at which a valid pattern of this weight turned up;
no exhaustive search was performed at that or any other period, and no certified minimum is
claimed. In particular we make no claim that $59/156$ is close to optimal: it is one interior
point of the recorded interval $[43/120,2/5]$, and the transfer-matrix and minimum-cycle-mean
machinery of \cite{Jiang}, \cite{BHP} and \cite{BDMP} may well settle $k=6$ outright. Only the
inequality $\dstar(S_6)\le59/156$ is asserted.

\section{Verification}

The accompanying program \texttt{verify.py} (Python~3.9+, standard library only, no external
data file, exact integer and \texttt{Fraction} arithmetic) reads the $52$ masks printed in
Section~2, regenerates the six row strings from them, and re-derives the row weights, the count
$118$ of $312$ and the density $59/156$ as an exact rational, and then verifies \eqref{eq:idc}:
no domination failures over the $312$ vertices of the fundamental domain and no separation
failures among the relevant pairs. Its recorded run is \texttt{verify.output.txt}, one
\texttt{PASS} line per check. The run also prints some checks corresponding to material deleted
from this note (the empty columns, the minimality of the period, the bracket arithmetic, and a
sweep of pairs at distance $3$ to $6$); those support nothing claimed here. It establishes the
upper bound of Theorem~\ref{thm:upper} and nothing else.

\begin{thebibliography}{9}

\bibitem{BDMP}
M.~Bouznif, J.~Darlay, J.~Moncel, and M.~Preissmann,
\emph{Exact values for three domination-like problems in circular and infinite grid graphs
of small height},
Discrete Math. Theor. Comput. Sci. \textbf{21} (2019), no.~3, paper~\#12.
\href{https://doi.org/10.23638/DMTCS-21-3-12}{doi:10.23638/DMTCS-21-3-12}.

\bibitem{BHP}
M.~Bouznif, F.~Havet, and M.~Preissmann,
\emph{Minimum-density identifying codes in square grids},
INRIA Research Report RR-8845 (hal-01259550v1); AAIM 2016,
Lecture Notes in Comput. Sci. \textbf{9778}, 77--88.
\href{https://doi.org/10.1007/978-3-319-41168-2_7}{doi:10.1007/978-3-319-41168-2\_7}.

\bibitem{DGM}
M.~Daniel, S.~Gravier, and J.~Moncel,
\emph{Identifying codes in some subgraphs of the square lattice},
Theoret. Comput. Sci. \textbf{319} (2004), 411--421.
\href{https://doi.org/10.1016/j.tcs.2004.02.007}{doi:10.1016/j.tcs.2004.02.007}.

\bibitem{Jiang}
M.~Jiang,
\emph{Periodicity of identifying codes in strips},
Inform. Process. Lett. \textbf{135} (2018), 77--84.
\href{https://doi.org/10.1016/j.ipl.2018.03.007}{doi:10.1016/j.ipl.2018.03.007};
\href{https://arxiv.org/abs/1607.03848v2}{arXiv:1607.03848v2}, which is the version
used here.

\bibitem{LHC}
A.~Lobstein, O.~Hudry, and I.~Charon,
\emph{Locating-Domination and Identification},
Chapter~6 of \emph{Topics in Domination in Graphs} (T.~W. Haynes, S.~T. Hedetniemi,
M.~A. Henning, eds.), Springer, 2020, 251--299.
\href{https://doi.org/10.1007/978-3-030-51117-3_8}{doi:10.1007/978-3-030-51117-3\_8}.
Theorem~30 is quoted from the open copy
\href{https://hal.science/hal-02916929/document}{hal-02916929v1}, PDF page~11
(manuscript page~10).

\end{thebibliography}

\end{document}
